% paper/sections/methodology.tex
% §4.1 GNN encoder + joint heads
% §4.2 LSTM sequence model
% §4.3 Training: imitation → REINFORCE + mixed expert
% §4.4 Calibrated DP baseline

\section{Method}
\label{sec:method}

Our model — Rev-GNN-LSTM — consists of a GraphSAGE encoder,
an LSTM episode-history module, and two joint output heads that simultaneously
select the next buyer and set a personalized discount.

\subsection{Node Feature Representation}
\label{sec:features}

Each node $i \in \calV$ at step $t$ is represented by a 21-dimensional feature vector
$\phi_i^{(t)} \in \R^{21}$ composed of three groups.

\emph{Static structural features} (dims 1–10, identical to~\citet{anon_wsdm27}):
degree rank, clustering coefficient, betweenness centrality, PageRank,
$k$-core number, eigenvector centrality, triangle count, closeness centrality,
eccentricity, and average neighbor degree.

\emph{Dynamic state features} (dims 11–16):
seed flag $\mathbf{1}[i \in \St]$, round ratio $t/n$,
fraction of $i$'s neighbors already in $\St$,
$\log(1 + \deg(i))$, clustering coefficient (repeated as diffusion signal),
and group flag (set to 0 for the revenue task).

\emph{Pricing-specific features} (dims 17–21, new):
normalized influence $x_i(\St)$,
estimated valuation $\hat{\val}_i(\St)$ from 200 Monte Carlo samples,
an offered-and-rejected flag,
steps-remaining ratio $(n - t)/n$,
and current budget ratio $B_t / B_0$ (dims 17–21; dim 21 is $0$ for the
unconstrained problem).

\subsection{GNN Encoder}
\label{sec:gnn}

A two-layer GraphSAGE~\citep{hamilton2017inductive} encoder with residual
connections and LayerNorm maps each node's feature vector to a 64-dimensional
embedding:
\begin{align}
  h_i^{(0)} &= \mathrm{ReLU}(\mathrm{LN}(W_{\mathrm{proj}}\,\phi_i^{(t)})),\\
  h_i^{(l)} &= \mathrm{ReLU}\!\left(\mathrm{LN}\!\left(h_i^{(l-1)}
    + W_{\mathrm{self}} h_i^{(l-1)}
    + W_{\mathrm{neigh}} \tfrac{1}{|\calN(i)|}\sum_{j\in\calN(i)} h_j^{(l-1)}
  \right)\!\right).
\end{align}
All weight matrices are shared across graphs and steps.

\subsection{LSTM Episode-History Module}
\label{sec:lstm}

A stateless GNN sees only the current step's graph embedding; it cannot reason
about pricing signals revealed by past offers, such as a network that
consistently rejects discounts below 40\%.
We address this with an LSTM over the episode trajectory.

At each step $t$, the LSTM receives a token
\[
  \tau_t = \bigl[\bar{h}^{(t)}\;\|\; d_{t-1}\;\|\; a_{t-1}\;\|\; r_{t-1}\bigr]
  \in \R^{67},
\]
where $\bar{h}^{(t)} = \frac{1}{|\calV|}\sum_i h_i^{(2)}$ is the mean-pooled
graph embedding, $d_{t-1}$ is the last discount offered, $a_{t-1} \in \{0,1\}$
is the acceptance indicator, and $r_{t-1}$ is the revenue collected.
The LSTM hidden state $c_t \in \R^{64}$ summarizes the full history.

\subsection{Joint Policy Heads}
\label{sec:heads}

For each unvisited node $i \notin \St$ the full context embedding is
\[
  z_i^{(t)} = \bigl[h_i^{(2)}\;\|\; c_t\bigr] \in \R^{128}.
\]
Two heads share this representation:

\emph{Scoring head}: $s_i = \mathrm{Linear}_{32\to1}(\mathrm{ReLU}(\mathrm{Linear}_{128\to32}(z_i^{(t)})))$.
The agent selects $v^* = \argmax_{i \notin \St} s_i$.

\emph{Pricing head}: a two-layer MLP that outputs the parameters
$(\alpha, \beta)$ of a Beta distribution over the discount $d_t \in [0,1]$:
\[
  (\alpha, \beta) = \mathrm{Softplus}\!\left(\mathrm{Linear}_{32\to2}(\mathrm{ReLU}(\mathrm{Linear}_{128\to32}(z_{v^*}^{(t)})))\right) + 1.
\]
During training the discount is sampled $d_t \sim \mathrm{Beta}(\alpha,\beta)$;
during evaluation we use the mean $d_t = \alpha/(\alpha+\beta)$.

\subsection{Training: Imitation then REINFORCE}
\label{sec:training}

Training proceeds in two phases, following the mixed-expert
schedule of~\citet{anon_wsdm27}.

\paragraph{Phase 1 — Imitation learning (300 epochs).}
We generate expert trajectories using the Greedy-Discount algorithm of
\citet{babaei2013revenue}.
For each expert trajectory we minimize a cross-entropy loss on the selection
head and a Beta-distribution negative log-likelihood on the pricing head.
Training graphs: 5 Forest-Fire graphs with $n \in \{200, 260, 320, 380, 440\}$,
drawn with $p=0.37$, $p_b=0.32$.
After Phase 1 the cross-entropy loss decreases from $4.96$ to $2.68$.

\paragraph{Phase 2 — REINFORCE (200 epochs).}
We fine-tune the full policy (selection $+$ pricing) with REINFORCE.
To stabilize training across the wide range of episode lengths and budget values,
we group episodes by budget size into five buckets
(corresponding to $k \in \{1\text{–}2, 3\text{–}5, 6\text{–}10, 11\text{–}20, 21\text{–}40\}$)
and maintain a per-bucket Welford running mean and standard deviation for
advantage normalization.
The entropy bonus coefficient is $\lambda_H = 0.01$ and the learning rate is
$5 \times 10^{-5}$.
The best checkpoint is selected by the maximum of the minimum per-bucket
mean reward.

A single policy trained across all budgets must reconcile opposing
strategies: capital-binding episodes reward conservative self-funding,
budget-slack episodes reward aggressive seeding.
To remove this compromise we additionally train a large-$k$ specialist on
$k \in [16, 40]$ only, warm-started from the unified model with the
identical recipe.
At deployment the checkpoint is selected by budget regime, mirroring the
per-$k$ composition of our Cal-DP baseline.
Both checkpoints are released.

\begin{figure}[t]
  \centering
  \includegraphics[width=0.95\columnwidth]{fig6_training_curve}
  \caption{Training progression of Rev-GNN-LSTM.
    \textbf{Left:} Phase~1 (imitation, 300 epochs) reduces cross-entropy from
    $4.96{\to}2.68$.
    \textbf{Right:} Phase~2 (REINFORCE, 200 epochs, best checkpoint at ep.\ 43)
    improves mean episode revenue on the training graphs.
    Per-bucket advantage normalization stabilizes training across the full $k$
    range.}
  \label{fig:training_curve}
\end{figure}

\begin{figure}[t]
  \centering
  \includegraphics[width=0.95\columnwidth]{fig5_discount_trajectory}
  \caption{Discount trajectory of Rev-GNN-LSTM over one episode (FF $n{=}1000$).
    Early buyers (high degree, low influence) receive deep discounts to
    bootstrap diffusion; late buyers (already highly influenced) are offered
    near-list prices as their valuations are elevated by peer purchases.
    This pattern emerges from RL fine-tuning, not from the imitation teacher.}
  \label{fig:discount_traj}
\end{figure}

\subsection{Calibrated Dynamic-Programming Baseline (Cal-DP)}
\label{sec:caldp}

For the budget-constrained setting we introduce a family of non-learned baselines
that calibrate classical DP to the replenishable-capital structure.

\paragraph{Cal-DP v2 (two-phase DP).}
Phase 1 uses a standard value-iteration DP over a discretized budget grid,
computing the expected revenue for each remaining budget level.
Phase 2 applies budget recalibration: the DP value function is rescaled to
account for the fact that accepted sales replenish the budget,
making the effective budget larger than its initial value.

\paragraph{Cal-DP v3 (tier-based allocation).}
Buyers are partitioned into tiers by estimated valuation.
A DP over tier assignment allocates budget to tiers in decreasing value order,
using a calibrated production-cost offset to account for the probabilistic cash
flow.
Cal-DP v3 dominates at small $k$ (capital-binding regime); Cal-DP v2 dominates
at large $k$ (budget-slack regime).

\paragraph{Composite Cal-DP.}
For the paper comparison we use $\max(\text{v2}, \text{v3})$ at each $k$,
taking the best algorithm per regime.
This composite dominates Cal-DP v1 (the standard DP) at every $k$ value tested
on the $n{=}1000$ Forest-Fire graph (see Table~\ref{tab:dp_family}).

\paragraph{No warm-start seeding.}
In both Cal-DP v2 and v3 there is no separate ``free-seed'' warm-start phase
(seed\_frac$=0$), ensuring a fair comparison with the learned policy which
faces the same budget constraint from step 1.
